% Results / Experiments (classified by basename + \section heading). NeurIPS norms: each subsection opens
% by naming the question it answers; the number is read out BEFORE the conclusion is drawn; the battery
% spans multiple models x datasets x baselines with ABLATIONS that isolate the mechanism. Use bold run-in
% experimental-setup paragraphs (\textbf{Models.} \textbf{Datasets.} \textbf{Baselines.} \textbf{Metrics.}
% \textbf{Implementation Details.}). Reproduce-then-compare. Report losses honestly with a mechanism.
% Efficiency claims: name the hardware and account for ALL costs; a real serving framework (SGLang/vLLM)
% for throughput. Tables: booktabs (no vertical rules), BEST value in bold, caption ABOVE the table.
\section{Experiments}
\label{sec:results}

\paragraph{Models.} [placeholder]
\paragraph{Datasets.} [placeholder]
\paragraph{Baselines.} [placeholder — the real competitors, reproduced at matched settings.]
\paragraph{Metrics.} [placeholder — the RIGHT metric for the claim.]
\paragraph{Implementation Details.} [placeholder — hyperparameters, seeds, hardware.]

\subsection{Main results}
\label{sec:results-main}
[Results placeholder — state the question, point to the table/figure, read out the numbers, then draw the
conclusion and explain any variance. Report loss cases with a mechanism.]

\subsection{Ablation study}
\label{sec:results-ablation}
[Ablation placeholder — remove each proposed component and show its contribution, isolating the mechanism
the paper credits for the gain.]
